% Versión en español del perfil DevOps/SRE Engineer Lead. Mantener cifras,
% alcance y las notas de disciplina de hechos alineados con lang/lead/en.tex
% (ver lang/lead/en.tex, cabecera de comentarios, y FACTS.md).
%
% En particular: la arquitectura hexagonal en Veradat fue REVISADA, no
% implementada; la planeación en Microsoft Project fue por sprint sobre
% actividades propias, sin gestión de dependencias ni hitos; el rol de
% ISO 27001 fue implementar y operar controles, no auditar ni certificar.

\renewcommand{\resumeaddress}{Ciudad de México, México \textbar\ UTC-6 (US Central)}
\resumejobtitle{DevOps/SRE Engineer Lead \textbar\ Remediación de deuda técnica, estándares y cumplimiento}
\resumesummary{DevOps/SRE con nueve años de experiencia, incluyendo trayectoria en desarrollo de software desde 2017 en banca y telecomunicaciones. Lidero la remediación de deuda técnica en 10 pipelines de CI/CD que atienden a tres equipos, e implemento y opero controles \textbf{ISO 27001} en Veradat y Banamex. Reviso decisiones de arquitectura del equipo --incluida la arquitectura hexagonal que desacopla los microservicios de PostgreSQL, Oracle, RabbitMQ, Kafka, IBM~MQ y SQS-- y genero la matriz de imágenes Docker que cada microservicio requiere. Acompaño a ingenieros mediante revisiones de código y diseño, y cubro guardia para incidentes en producción.}
\makeheader

\section*{Competencias Técnicas}
\skillrow{Plataformas y orquestación}{Kubernetes, OpenShift, Docker, Helm}
\skillrow{Nube e infraestructura como código}{AWS (Route~53, IAM), Azure, Terraform, Ansible, Linux, Bash}
\skillrow{CI/CD y entrega}{Azure DevOps, Jenkins, Nexus, GitHub, Maven}
\skillrow{Confiabilidad y respuesta a incidentes}{Elastic Stack (ELK), Splunk, ServiceNow, sondas readiness/liveness, guardia de producción}
\skillrow{Cumplimiento y seguridad}{ISO 27001 (implementación y operación de controles), SonarQube Cloud, Snyk, Mend}
\skillrow{Arquitectura y prácticas}{Microservicios, revisión de arquitectura hexagonal, adaptadores de integración legada, diseño de APIs REST, Scrum, Microsoft Project}
\skillrow{Desarrollo}{Java, Spring Cloud, JPA, Python, Go, Kafka, RabbitMQ, IBM~MQ}

\section*{Experiencia Profesional}

\resumeentry{DevOps Engineer}{Veradat Smart Network S.A. de C.V.}{Ago 2024 -- Actualidad}{Ciudad de México}
\begin{itemize}
    \item Diseñé y mantuve \textbf{10 pipelines de CI/CD en Azure DevOps} para \textbf{3 equipos} (frontend, backend, QA), reduciendo el tiempo de despliegue del conjunto completo de 70 microservicios de 6 horas a 1 hora (\textasciitilde 83\% menos). Desarrollé y promoví \textbf{Helm charts como estándar de despliegue del equipo}, estandarizando y versionando el empaquetado de despliegues en los tres equipos.
    \item Revisé la \textbf{arquitectura hexagonal} del equipo, que desacopla cada microservicio de sus dependencias tecnológicas (PostgreSQL, Oracle, RabbitMQ, Kafka, IBM~MQ, SQS), y generé la \textbf{imagen Docker para cada combinación tecnológica} que cada microservicio requiere.
    \item \textbf{Implementé y operé controles ISO 27001} en gestión de secretos, control de acceso en nube y escaneo de dependencias, y remedié \textbf{vulnerabilidades críticas y altas} detectadas por SonarQube Cloud, Snyk y Mend mediante actualizaciones de paquetes del sistema operativo y dependencias Maven.
    \item \textbf{Acompaño} a ingenieros de los equipos de frontend, backend y QA mediante \textbf{revisiones de código y diseño}, decisiones técnicas e incorporación de nuevos miembros al equipo.
    \item Administro \textbf{9 clústeres de producción (8 Kubernetes, 1 OpenShift)}, con sondas readiness/liveness para controlar el tráfico durante actualizaciones progresivas y disparar la recuperación. Cubro \textbf{guardia para incidentes en producción}; implementé monitoreo y alertas con \textbf{Elastic Stack} en los 9 clústeres, reduciendo el \textbf{tiempo de detección de incidentes de 60 a 30 minutos} (50\% menos).
    \item Aprovisioné infraestructura en \textbf{AWS y Azure con Terraform} --máquinas virtuales, balanceadores de carga, gestión de secretos, Route~53 e IAM-- y automaticé la configuración con Ansible en servidores Linux, reduciendo el aprovisionamiento del conjunto de 10 servidores de 2 horas a 20 minutos (\textasciitilde 83\% menos).
    \item Planeé y prioricé actividades de infraestructura en \textbf{Microsoft Project} por sprint, en línea con la cadencia de entrega de los tres equipos.
    \item Desarrollé herramientas de automatización y tooling interno en \textbf{Python (\textasciitilde 2 años) y Go (1 año)} para operaciones de infraestructura y monitoreo, usando \textbf{Claude Code (\textasciitilde 1 año)} como asistente de desarrollo.
\end{itemize}

\resumeentry{DevOps Engineer Fullstack}{IDS Comercial S.A. de C.V. \textbar\ Cliente: Banamex}{Ene 2021 -- Ago 2024}{México, remoto}
\begin{itemize}
    \item \textbf{Estandemicé la plataforma con un POM padre de Maven multimódulo} y gestión compartida de dependencias, garantizando versionado consistente entre servicios y reduciendo el esfuerzo de actualización -- una reducción directa de deuda de mantenimiento.
    \item Desarrollé y mantuve \textbf{10 microservicios Java/Spring Cloud en producción} para flujos de trabajo de ejecutivos bancarios que atendían aproximadamente \textbf{1,000 transacciones diarias por sucursal}.
    \item Implementé y operé controles \textbf{ISO 27001} en este entorno bancario, y planeé y prioricé actividades de infraestructura en \textbf{Microsoft Project} por sprint.
    \item Construí pipelines de CI/CD con \textbf{Azure DevOps y Jenkins} para automatizar compilaciones y despliegues de la aplicación bancaria.
    \item Atendí y prioricé incidentes de producción y QA bajo un \textbf{SLA de 1 día en ServiceNow}, con un tiempo promedio de resolución de 2 a 3 horas y casos complejos de hasta una semana; monitoreé producción con Splunk.
\end{itemize}

\resumeentry{DevOps Engineer de Software Fullstack}{Proyecto continuo para Telcel, mediante GLOBAL HITSS y Everis}{Jul 2018 -- Dic 2020}{Ciudad de México}
\textit{GLOBAL HITSS: jul 2018 -- sep 2019; Everis: sep 2019 -- dic 2020.}\par\vspace{3pt}
\begin{itemize}
    \item \textbf{Construí una capa de integración con Java, Spring Cloud y JPA} que tradujo protocolos de sistemas legados de telecomunicaciones a APIs REST, conectando sistemas existentes con servicios modernos -- una capa anticorrupción entre el legado y el estado actual.
    \item \textbf{Implementé pipelines de CI/CD y contenericé el servicio con Docker sobre Kubernetes.} Automaticé procesos operativos con Bash y gestioné repositorios de código y artefactos con GitHub y Nexus.
\end{itemize}

\section*{Proyecto Destacado}
\resumeentry{Sistema de Gestión de Cumplimiento de la Ley de Subcontratación}{GDP administración S de RL de CV (Grupo de la Parra)}{Proyecto freelance, liderazgo técnico}{Ciudad de México}
\begin{itemize}
    \item \textbf{Lideré el diseño técnico y la arquitectura end-to-end} de una plataforma en producción para gestionar el cumplimiento de la \textbf{Ley Federal del Trabajo en materia de subcontratación} sobre una fuerza laboral subcontratada de \textbf{5,000 empleados} y múltiples empresas proveedoras y receptoras. \textbf{Acompañé al equipo del cliente} durante el diseño.
    \item Construí un \textbf{motor de auditoría de cumplimiento} que valida reglas regulatorias (existencia de CFDI por empleado vs. fechas de pago, pagos de IMSS/INFONAVIT/retenciones vs. calendario, alta ante la Secretaría del Trabajo) y emite \textbf{alertas automatizadas} por vencimientos de contratos y plazos de entrega de información.
    \item Entregué un \textbf{panel de BI de cumplimiento} con reportes regulatorios (\textbf{SAT, IMSS, INFONAVIT, STPS}) y un dashboard con KPIs (documentos pendientes, días para entrega, contratos por proveedor/cliente), que permitió al área legal y de antifraude operar de forma proactiva.
\end{itemize}

\section*{Educación}
\resumeentry{IPN -- Escuela Superior de Cómputo (ESCOM)}{Ingeniería en Sistemas Computacionales}{2014 -- 2018}{Ciudad de México}

\section*{Certificación e Idiomas}
\skillrow{Certificación}{Java SE -- Oracle, Certificado Asociado (2013). \href{https://www.youracclaim.com/badges/583548b1-06fa-4dbf-889d-419e6c6c2ec1/public_url}{Credencial}; Scrum Fundamentals -- SCRUMstudy (2018).}
\skillrow{Idiomas}{Español nativo; inglés B2 (intermedio alto)}

\section*{Reconocimientos}
\begin{itemize}
    \item Participante, Google Code Jam (2019); equipo de selección de etapa de ACM-ICPC (2015--2017).
\end{itemize}
